\documentclass[10pt]{mypackage}

\usepackage[uprightgreeks,noDcommand,nofligatures]{kpfonts}
\renewcommand{\coloneq}{\coloneqq}
\renewcommand{\eqcolon}{\eqqcolon}
%\usepackage{dsfont}
%\renewcommand{\mathbb}{\mathds}
%\usepackage{homework}
\usepackage{notes}

\usepackage[ backend=bibtex, style = alphabetic, sorting=ynt ]{biblatex}
\addbibresource{../all_references.bib}

\usepackage{parskip}
\frenchspacing

\fancyhf{}
\fancyhead[R]{Avinash Iyer}
\fancyhead[L]{Regular Conditional Measures, Disintegration, and The Ergodic Decomposition}
\fancyfoot[C]{\thepage}

\setcounter{secnumdepth}{0}

\begin{document}
\RaggedRight
\section{Regular Conditional Measures}%
\subsection{Definitions}%
Suppose $\mu$ is a measure on $\left( X,\mathcal{A} \right)$. Let $\mathcal{B}$ be a sub-$\sigma$-algebra of $ \mathcal{A} $. We may view $B = \sigma\left( \pi \right)$ for some measurable map $\pi\colon \left( X,\mathcal{A} \right)\to \left( Y,\mathcal{E} \right)$.

For a positive measure $\mu$, and every $A\in \mathcal{A}$, there exists a $\mathcal{B}$-measurable function $\mu\left( A,\cdot \right)$ such that
\begin{align*}
  \mu\left( A\cap B \right) &= \int_{B}^{} \mu\left( A,x \right)\:d\mu(x)
\end{align*}
for all $B\in \mathcal{B}$. Specifically, this is the Radon--Nikodym derivative of the measure $\mu\left( A\cap \cdot \right)\ll \mu$. Since the map $x\mapsto \mu\left( A,x \right)$ is $\mathcal{B}$-measurable, there exists some $ \mathcal{E} $-measurable function $y\mapsto \mu^y(A)$ such that $\mu\left( A,x \right) = \mu^{\pi(x)}\left( A \right)$. If we let $\nu = \pi_{\ast}\mu$, then
\begin{align*}
  \mu\left( A\cap \pi^{-1}\left( E \right) \right) &= \int_{E}^{} \mu^y(A)\:d\nu(y)
\end{align*}
for all $E\in \mathcal{E}$. In particular, we may let $E = Y$, and we have
\begin{align*}
  \mu(A) &= \int_{Y}^{} \mu^y(A)\:d\nu(y).
\end{align*}
We may view $\mu^y$ as a measure on $\pi^{-1}\left( \set{y} \right)$ for every $y\in Y$; if so, then we can view this as a generalized Fubini's Theorem. That is, to compute the measure of $A\in \mathcal{A}$, we compute the conditional measures of $A$ on the level sets $\pi^{-1}\left( \set{y} \right)$ and integrate in $y$ with respect to the measure $\nu$.

This measure may not be countably additive; to find such measures that are, we will require some additional assumptions.
\begin{definition}
  Let $\mathcal{A}$ be a $\sigma$-algebra on $X$, $\mathcal{B}$ a sub-$\sigma$-algebra, and $\mu$ a measure on $ \mathcal{A} $. We say
  \begin{align*}
    \mu^{\mathcal{B}}\left( \cdot,\cdot \right) \colon \mathcal{A}\times X &\to \R
  \end{align*}
  is a \textit{regular conditional measure} (or system of regular conditional measures) on $ \mathcal{A} $ with respect to $ \mathcal{B} $ if
  \begin{enumerate}[(i)]
    \item for every $x\in X$, the function $A\mapsto \mu^{\mathcal{B}}\left( A,x \right)$ is a measure on $ \mathcal{A} $;
    \item for every $A\in \mathcal{A}$, the function $x\mapsto \mu^{\mathcal{B}}\left( A,x \right)$ is measurable with respect to $ \mathcal{B} $ and $ \left\vert \mu \right\vert $-integrable;
    \item for all $A\in \mathcal{A}$ and $B\in \mathcal{B}$, we have
      \begin{align*}
        \mu\left( A\cap B \right) &= \int_{B}^{} \mu^{\mathcal{B}}\left( A,x \right)\:d\left\vert \mu \right\vert(x).
      \end{align*}
  \end{enumerate}
\end{definition}
\begin{definition}
  If $\pi\colon X\to Y$ is a $\mathcal{A}$-$\mathcal{E}$-measurable map, then a system of regular conditional measures is a function $\left( A,y \right)\mapsto \mu^y(A)$ on $ \mathcal{A}\times Y $ such that for every fixed $y$, it is a measure on $ \mathcal{A} $, for every fixed $A\in \mathcal{A}$, it is measurable and $\pi_{\ast}\left\vert \mu \right\vert$-integrable, and satisfies for all $A\in \mathcal{A}$ and $E\in \mathcal{E}$ that
  \begin{align*}
    \mu\left( A\cap \pi^{-1}\left( E \right) \right) &= \int_{E}^{} \mu^y(A)\:d\pi_{\ast}\left\vert \mu \right\vert.
  \end{align*}
  If $\pi^{-1}\left( \set{y} \right)\in \mathcal{A}$ for $\pi_{\ast}\left\vert \mu \right\vert$-a.e. $y\in Y$, and $\mu^y$ is concentrated on $\pi^{-1}\left( \set{y} \right)$, then we call the family \textit{proper} conditional measures.
\end{definition}
We may generalize this further by considering a sub-$\sigma$-algebra $\mathcal{A}_0$ (that may not contain $\mathcal{B}$) and mandating the equality
\begin{align*}
  \mu\left( A\cap B \right) &= \int_{B}^{} \mu^{\mathcal{B}}_{\mathcal{A}_0}\left( A,x \right) \:d\left\vert \mu \right\vert(x)
\end{align*}
hold for all $A\in \mathcal{A}_0$ and $B\in \mathcal{B}$.
\subsection{Existence and Uniqueness on a Single Space}%
We note in the following statement that in most settings (i.e., standard probability spaces), we will be working with countable generation.
\begin{proposition}[Uniqueness]
  Suppose $\mathcal{A}$ is countably generated. Then, given any two regular conditional measures $\mu_1^{\mathcal{B}}$ and $\mu_2^{\mathcal{B}}$ on $ \mathcal{A} $, there is $Z\in \mathcal{B}$ with $\left\vert \mu \right\vert\left( Z \right) = 0$ such that $\mu_1^{\mathcal{B}}\left( A,x \right) = \mu_2^{\mathcal{B}}\left( A,x \right)$ for all $A\in \mathcal{A}$ and all $x\in X\setminus Z$.

  Similarly, the measures $\mu^{\mathcal{B}}_{\mathcal{A}_0}\left( \cdot,x \right)$ on $ \mathcal{A}_0 $ are essentially unique if $\mathcal{A}_0$ is countably generated (even if $ \mathcal{A} $ may not be countably generated).

  If $\mu$ is a probability measure, then $\mu^{\mathcal{B}}\left( \cdot,x \right)$ is a probability measure for $\mu$-a.e. $x$.
\end{proposition}
\begin{proof}
  Let $ \mathcal{R} = \set{A_n : n\in \N} $ be a countable algebra generating $ \mathcal{A} $. Since for all $B\in \mathcal{B}$, we have that
  \begin{align*}
    \mu\left( A_n\cap B \right) &= \int_{B}^{} \mu_{1}^{\mathcal{B}}\left( A,x \right)\:d\left\vert \mu \right\vert(x)\\
                                &= \int_{B}^{} \mu_2^{\mathcal{B}}\left( A,x \right)\:d\left\vert \mu \right\vert(x),
  \end{align*}
  it follows by the essential uniqueness of Radon--Nikodym derivatives that there is some $Z_n\in \mathcal{B}$ such that $\mu_1^{\mathcal{B}}\left( A_n,x \right) = \mu_2^{\mathcal{B}}\left( A_n,x \right)$ for all $x\in X\setminus Z_n$. If we take $Z = \bigcup_{n=1}^{\infty}Z_n$, then we obtain the desired result.

  If $\mu$ is a probability measure, then for all $n\in \N$, the function $\mu^{\mathcal{B}}\left( A_n,x \right)$ is positive $\mu$-a.e. since its integral over every $B\in \mathcal{B}$ is positive. Similarly, $\mu^{\mathcal{B}}\left( X,x \right) = 1$ for $\mu$-a.e. $x$, so for $\mu$-a.e. $x$, the measure $\mu^{\mathcal{B}}\left( \cdot,x \right)$ is positive on $ \mathcal{R} $ and endows $X$ with measure $1$, so $\mu^{\mathcal{B}}\left( \cdot,x \right)$ is a probability measure for all such $x$.
\end{proof}
\begin{definition}
  A family of subsets $\mathcal{K}$ is called a \textit{compact approximating class} for $ \mu $ if
  \begin{itemize}
    \item it is a compact class, in the sense that if $\bigcap_{n=1}^{\infty}K_n = \emptyset$ for some family of $K_n\in \mathcal{K}$, then there is some $N$ such that $\bigcap_{n=1}^{N}K_n = \emptyset$;
    \item for every $A\in \mathcal{A}$ and every $\ve > 0$, there exists $K_{\ve}\in \mathcal{K}$ and $A_{\ve}\in \mathcal{A}$ such that $A_{\ve}\subseteq K_{\ve}\subseteq A$ and $\mu\left( A\setminus A_{\ve} \right) < \ve$
  \end{itemize}
\end{definition}
\begin{theorem}[Existence]\hfill
  \begin{enumerate}[(i)]
    \item Suppose the $\sigma$-algebra $\mathcal{A}$ is countably generated and $\mu$ has a compact approximating class in $ \mathcal{A} $. Then, for every sub-$\sigma$-algebra $ \mathcal{B} $, there exists a regular conditional measure $\mu^{\mathcal{B}}$ for $ \mathcal{A} $.
    \item Suppose $\mathcal{A}_0$ is a sub-$\sigma$-algebra of $ \mathcal{A} $ that is generated by a countable algebra $ \mathcal{U} $, and there is a compact class $\mathcal{K}\subseteq \mathcal{A}$ such that for all $A\in \mathcal{U}$ and for all $\ve > 0$, there exists $K_{\ve}\in \mathcal{K}$ and $A_{\ve}\in \mathcal{A}$ such that $A_{\ve}\subseteq K_{\ve}\subseteq A$ with $\left\vert \mu \right\vert\left( A\setminus A_{\ve} \right) < \ve$. Then, for every sub-$\sigma$-algebra $\mathcal{B}\subseteq \mathcal{A}$, there exists a regular conditional measure $\mu^{\mathcal{B}}_{\mathcal{A}_0}$ on $\mathcal{A}_0$.

      Moreover, for every $\mathcal{A}_0$-measurable $\mu$-integrable function $f$, one has
      \begin{align*}
        \int_{X}^{} f\:d\mu &= \int_{X}^{} \int_{X}^{} f(y)\:d\mu^{\mathcal{B}}_{\mathcal{A}_0}\left( y,x \right)\:d\left\vert \mu \right\vert\left( x \right).
      \end{align*}
  \end{enumerate}
\end{theorem}
\begin{proof}
  We will prove the second assertion, and start by assuming that $\mu$ is a probability measure. Let $ \mathcal{U} = \set{A_n : n\in \N}$ be a countable generating family for $ \mathcal{A}_0 $. For every $n\in \N$, we find $C_{n,k}\in \mathcal{K}$ and $A_{n,k}\in \mathcal{A}$ such that $A_{n,k}\subseteq C_{n,k}\subseteq A_n$ and $\mu\left( A_n\setminus A_{n,k} \right) < 1/k$.

  The sets $A_{n,k}$ along with $A_n$ generate a countable algebra $\mathcal{U}_0\subseteq \mathcal{A}$. For each $A\in \mathcal{U}_0$, we have $\mathcal{B}$-measurable Radon--Nikodym derivatives $x\mapsto p_0\left( A,x \right)$ with $p_0\left( X,x \right) = 1$, $p_0\left( \emptyset,x \right) = 0$ for all $x$, and
  \begin{align*}
    \mu\left( A\cap B \right) &= \int_{B}^{} p_0\left( A,x \right)\:d\mu(x).
  \end{align*}
  We claim that there exists a measure zero subset $N_0\in \mathcal{B}$ such that for all $x\in X\setminus N_0$, the function $A\mapsto p_0\left( A,x \right)$ is additive on $\mathcal{U}_0$. It follows by the integral expression for $\mu\left( A\cap B \right)$ that $p_0\left( A\cup B,x \right) = p_0\left( A,x \right) + p_0\left( B,x \right)$ $\mu$-a.e. whenever $A,B\in \mathcal{U}_0$ with $A\cap B = \emptyset$. Since the set of all such pairs $\left( A,B \right)$ in $\mathcal{U}_0$ is countable, the desired result follows.

  Next, we show that
  \begin{align*}
    p_0\left( A_n,x \right) &= \sup_{k}p_0\left( A_{n,k},x \right).
  \end{align*}
  Denote the right-hand side as $q_n(x)$. The function $q_n(x)$ is $\mathcal{B}$-measurable. Since $A_{n,k}\subseteq A_n$, we have that there are measure zero subsets $N_{n,k}\in \mathcal{B}$ such that $p_0\left( A_{n,k},x \right)\leq p_0\left( A_n,x \right)$ for all $x\in N_{n,k}$, so that
  \begin{align*}
    q_n(x) &\leq p_0\left( A_n,x \right)
  \end{align*}
  for all $x\notin N_0\coloneq \bigcup_{n,k=1}^{\infty}N_{n,k}$. Yet, since $p_0\left( A_{n,k},x \right)\leq q_n(x)$, we have that
  \begin{align*}
    \mu\left( A_{n,k} \right) &= \int_{X}^{} p_0\left( A_{n,k},x \right)\:d\mu(x)\\
                              &\leq \int_{X}^{} q_n(x)\:d\mu(x).
  \end{align*}
  Since $\mu\left( \bigcup_{n,k=1}^{\infty}N_{n,k} \right) = 0$, we have
  \begin{align*}
    \sup_{k}\mu\left( A_{n,k} \right) &\leq \int_{X}^{} q_n(x)\:d\mu(x)\\
                                      &\leq \int_{X}^{} p_0\left( A_n,x \right)\:d\mu(x)\\
                                      &= \mu\left( A_n \right).
  \end{align*}
  Since $\sup_{k}\mu\left( A_{n,k} \right) = \mu\left( A_n \right)$, we have that $ q_n(x) = p_0\left( A_n,x \right) $ everywhere except for some measure-zero $N_1\in \mathcal{B}$.

  Thus, for all $x\notin N = N_0\cup N_1$, we have that $p_0\left( \cdot,x \right)$ is an additive set function on the algebra $\mathcal{U}_0$ such that the compact class $\mathcal{K}$ approximates $p_0\left( \cdot,x \right)$ on $\mathcal{U}$ with respect to $\mathcal{U}_0$. So, this set function is countably additive on $\mathcal{U}$, meaning it extends to a unique countably additive measure on $ \mathcal{A}_0 $ by the Carathéodory Extension Theorem, which we define as $\mu\left( \cdot,x \right)\coloneq \mu^{\mathcal{B}}_{\mathcal{A}_0}\left( \cdot,x \right)$.

  We will seek to verify that this is our desired system of conditional measures. If $A = A_n$, then t he function $x\mapsto \mu\left( A,x \right)$ is measurable with respect to $\mathcal{B}$. Since the class of all sets $A\in \sigma\left( \mathcal{U} \right)$ for which this is true is monotone, it follows that it coincides with $\sigma\left( \mathcal{U} \right)$. If $B\in \mathcal{B}$ and $A = A_n$, then by construction one has
  \begin{align*}
    \mu\left( A\cap B \right) &= \int_{B}^{} \mu\left( A,x \right)\:d\mu(x).
  \end{align*}
  For a fixed $B\in \mathcal{B}$, the class $\mathcal{E}$ of all sets $A\in \sigma\left( \mathcal{U} \right)$ such that the above holds is also monotone, as if $E_j\nearrow E$, then $\lim_{j\to\infty}\mu\left( E_j,x \right) = \mu\left( E,x \right)$, whence by dominated convergence we have that $E\in \mathcal{E}$. Thus, $\mathcal{E} = \sigma\left( \mathcal{U} \right)$. Since $\sigma\left( \mathcal{U} \right) = \mathcal{A}_0$, it follows that
  \begin{align*}
    \mu\left( A\cap B \right) &= \int_{B}^{} \mu\left( A,x \right)\:d\mu(x)
  \end{align*}
  for all $A\in \mathcal{A}_0$ and all $B\in \mathcal{B}$. Bootstrapping then yields
  \begin{align*}
    \int_{X}^{} f\:d\mu &= \int_{X}^{} \int_{X}^{} f(y)\:d\mu(y,x)\:d\mu(x)
  \end{align*}
  for all $\mathcal{A}_0$-measurable $\mu$-integrable $f$.

  Now, if $\mu$ is a positive (but non-probabilistic) measure on $X$, then the conditional measures for $\frac{\mu}{\norm{\mu}}$ are conditional measures for $\mu$ as well. Finally, if $\mu$ is a signed measure, then we construct them as the difference of conditional measures $\mu^{+}$ and $\mu^{-}$, where the conditional measures for $\mu^{+}$ are concentrated on $X^{+}$ and the conditional measures on $\mu^{-}$ are concentrated on $X^{-}$, extending by $0$ where necessary.
\end{proof}
\subsection{Existence and Uniqueness Generated by a Measurable Map}%
Given a measurable space $\left( Y,\mathcal{E} \right)$ and a measure $\mu$, we will let $\mathcal{E}_{\mu}$ be the completion of $\mathcal{E}$ with respect to the measure $\mu$.
\begin{theorem}
  Let $\mu$ be a (signed) measure on a measurable space $\left( X,\mathcal{A} \right)$, $\left( Y,\mathcal{E} \right)$ a measurable space, and $\pi\colon \left( X,\mathcal{A} \right)\to \left( Y,\mathcal{E} \right)$ an $\mathcal{A}_{\mu}$-$\mathcal{E}$-measurable map. Suppose $\pi(X)\in \mathcal{E}_{\pi_{\ast}\left\vert \mu \right\vert}$, and suppose $\mathcal{A}$ is countably generated with a compact approximating class.

  Then, there exist regular conditional measures $\mu^y$ generated by $\pi$ on $\mathcal{A}$; these regular conditional measures are probability measures if $\mu$ is positive.

  More generally, $\pi$ generates regular conditional measures $\mu^y_{\mathcal{A}_0}$ for each $y\in Y$ on every countably generated sub-$\sigma$-algebra $\mathcal{A}_0\subseteq \mathcal{A}$ on on which $\mu$ has a compact approximating class.
\end{theorem}
\begin{proof}
  Let $\mathcal{B} = \sigma\left( \pi \right) = \pi^{-1}\left( \mathcal{E} \right)$. Reducing to the case with $\mathcal{A}_0 = \mathcal{A}$ and $\mathcal{A} = \mathcal{A}_{\mu}$, it follows that there is a conditional measure $\mu^{\mathcal{B}}\left( A,x \right)$ on $\mathcal{A}$ such that for every $A\in \mathcal{A}$, there is a function $\eta\left( A,y \right)$ such that
  \begin{align*}
    \mu^{\mathcal{B}}\left( A,x \right) &= \eta\left( A,\pi(x) \right).
  \end{align*}
  By hypothesis, there is $Y_0\in \mathcal{E}$ with $Y_0\subseteq \pi(X)$ and $\pi_{\ast}\left\vert \mu \right\vert\left( Y\setminus Y_0 \right) = 0$.

  So, for every $y\in Y_0$, $\eta\left( A,y \right)$ is a measure as a function of $A$. We will define $\mu^y = \eta\left( \cdot,y \right)$, and whenever $y\notin Y_0$, then $\mu^y = \mu$. For every $A\in \mathcal{A}$, we have that the function $\mu^y\left( A \right)$ is $\mathcal{E}$-measurable since $Y\setminus Y_0\in \mathcal{E}$. 
\end{proof}
\begin{remark}
  If $\mathcal{A}_0= \xi^{-1}\left( \mathcal{F} \right)$ is a produced $\sigma$-algebra by some (other) map $\xi\colon X\to \left( Z,\mathcal{F} \right)$, with $\nu = \xi_{\ast}\left\vert \mu \right\vert$ and $\xi(X)\in \mathcal{F}_{\nu}$, then if $\mathcal{F}$ is countably generated and the measure $\nu$ on $\mathcal{F}$ has a compact approximating class, then there are regular conditional measures $\mu^y_{\mathcal{A}_0}$ generated by $\pi$ on $\mathcal{A}_0$.
\end{remark}
\begin{corollary}
  Suppose that in the conditions of the above theorem, the $\sigma$-algebra $\mathcal{E}$ is countably generated and contains all singletons. Then, there exist regular conditional measures $\mu^y$ on the $\sigma$-algebra $\mathcal{A}' = \sigma\left( \mathcal{A},\pi^{-1}\left( \mathcal{E} \right) \right)$ such that for $\pi_{\ast}\left\vert \mu \right\vert$-a.e. $y$, the measure $\mu^y$ is concentrated on $\pi^{-1}\left( \set{y} \right)$. If $\pi$ admits an a.e. equal $\mathcal{A}$-$\mathcal{E}$-measurable version $ \widetilde{\pi} $ with $\widetilde{\pi}\left( \mathcal{A} \right)\subseteq \mathcal{E}_{\pi_{\ast}\left\vert \mu \right\vert}$, then $\pi^{-1}\left( \set{y} \right)\in \mathcal{A}_{\mu^{y}}$ for $\pi_{\ast}\left\vert \mu \right\vert$-a.e. $y$.
\end{corollary}
\begin{proof}
  We only need to consider the case of probability measures. Let $\nu = \pi_{\ast}\mu$. Then, by assumption, we have $\pi^{-1}\left( \set{y} \right)\in \mathcal{A}'$. There exists a countable algebra $\mathcal{E}_0 = \set{E_n : n\in \N}$ generating $\mathcal{E}$. It follows that $\mathcal{A}'$ is countably generated as well.

  By the above theorem, there exist regular conditional measures $\mu^y$ on $\mathcal{A}'$. Fix $E_n\in \mathcal{E}_0$. Then, for every $E\in \mathcal{E}$, we have
  \begin{align*}
    \int_{E}^{} \mu^y\left( \pi^{-1}\left( E_n \right) \right)\:d\nu(y) &= \mu\left( \pi^{-1}\left( E \right)\cap \pi^{-1}\left( E_n \right) \right)\\
                                                                        &= \mu\left( \pi^{-1}\left( E\cap E_n \right) \right)\\
                                                                        &= \nu\left( E\cap E_n \right)\\
                                                                        &= \int_{E}^{} 1_{E_n}(y)\:d\nu(y).
  \end{align*}
  So, $\mu^y\left( \pi^{-1}\left( E_n \right) \right) = 1_{E_n}(y)$ $\nu$-a.e. So, there exists a $\nu$-conull set $Y_0$ such that $\mu^y\left( \pi^{-1}\left( E_n \right) \right) = 1_{E_n}(y)$ for all $y\in Y_0$ and all $n$. So, bootstrapping, it follows that $\mu^y\left( \pi^{-1}\left( E \right) \right) = 1_E(y)$ for all $y\in Y_0$ and all $E\in \mathcal{E}$, since these are measures as functions of $E$ that coincide on the $\pi$-system $\mathcal{E}_0$. So, $\mu^y\left( \pi^{-1}\left( \set{y} \right) \right) = 1$.

  If there is a version $ \widetilde{\pi} $ that is $\mathcal{A}$-$\mathcal{E}$-measurable, then there exists some conull $X_0\subseteq X$ on which $ \pi $ coincides with $ \widetilde{\pi} $, so $Y_0 \coloneq \pi\left( X_0 \right) = \widetilde{\pi}\left( X_0 \right)$ is contained in $\mathcal{E}_{\nu}$. The measure $\mu$ on $X_0$ then admits a system of regular conditional measures $\mu^y$ for each $y\in Y_0$, defined on $\mathcal{A}_{X_0} = \mathcal{A}\cap X_0$, where $\mu^y\left( X_0\cap \pi^{-1}\left( \set{y} \right) \right) = 1$ for $\nu$-a.e. $y\in Y_0$. Extending these measures $\mu^y$ to $ \mathcal{A} $ by setting $\mu^y\left( X\setminus X_0 \right)$, we have that $X_0\cap \pi^{-1}\left( \set{y} \right)$ is in $\mathcal{A}$ and is contained in $\pi^{-1}\left( \set{y} \right)$.
\end{proof}
The last case we concern ourselves with is when $\Omega = X\times Y$ is a product space on measurable spaces $\left( X,\mathcal{A}_X \right)$ and $\left( Y,\mathcal{A}_Y \right)$ with $\sigma$-algebra $\mathcal{A} = \mathcal{A}_X\otimes \mathcal{A}_Y$. Let $\mathcal{B}_X$ and $\mathcal{B}_Y$ be the sub-$\sigma$-algebras in $\mathcal{A}$ formed by the sets $A\times Y$ and $X\times B$ with $A\in \mathcal{A}_X$ and $B\in \mathcal{A}_Y$ respectively. Define $\mu_X$ and $\mu_Y$ to be the pushforwards of $\mu$ on $X$ and $Y$, and let $\left\vert \mu \right\vert_X$ and $\left\vert \mu \right\vert_Y$ be the corresponding total variation measures. We will also define the sections
\begin{align*}
  A_y &= \set{x\in X : \left( x,y \right)\in A}
  A^y &= A\cap \left( X\times \set{y} \right)\\
      &= A_y\times \set{y}.
\end{align*}
\begin{theorem}
  Suppose $\mathcal{A}_X$ is countably generated and $\left\vert \mu \right\vert$ has a compact approximating class. Then, for every $y\in Y$, there is a measure $\mu\left( \cdot,y \right)$ on $\mathcal{A}_X$ such that the function $y\mapsto \mu\left( A,y \right) = \mu_y\left( A_y \right)$ is measurable with respect to $\left\vert \mu \right\vert_Y$ for every $A\in \mathcal{A}$. Moreover, for every $B\in \mathcal{A}_Y$, we have
  \begin{align*}
    \mu\left( A\cap \left( X\times B \right) \right) &= \int_{B}^{} \mu\left( A,y \right)\:d\left\vert \mu \right\vert_Y(y)\\
                                                    &= \int_{B}^{} \mu_y\left( A_y \right)\:d\left\vert \mu \right\vert_Y(y).
  \end{align*}
  We have that $\mu\left( A,y \right) = \mu\left( A^y,ya \right)$ if $\mathcal{A}_Y$ contains the singletons.

  Furthermore, for every $\mathcal{A}$-measurable $\mu$-integrable function $f$, we have
  \begin{align*}
    \int_{\Omega}^{} f\left( x,y \right)\:d\mu\left( x,y \right) &= \int_{Y}^{} \int_{X}^{} f\left( x,y \right)\:d\mu_y(x)\:d\left\vert \mu \right\vert_Y(y).
  \end{align*}
\end{theorem}
\begin{proof}
  The sets in $\mathcal{B}_X$ are of the form $E\times Y$ with $E\in \mathcal{A}_X$. There exist probability measures $\mu^y_{\mathcal{B}_X}$ with $y\in Y$ on $\mathcal{B}_X$ such that for all $E\in \mathcal{A}_X$ and $B \in \mathcal{A}_Y$, we have
  \begin{align*}
    \mu\left( \left( E\cap Y \right)\cap \left( X\times B \right) \right) &= \int_{B}^{} \mu^y_{\mathcal{B}_X}\left( E\times Y \right)\:d\mu_Y(y),
  \end{align*}
  where the integrand is $\mathcal{A}_Y$-measurable. Define probability measures $\mu_y$ on $\mathcal{A}_X$ by taking $\mu_y(E) = \mu^y_{\mathcal{B}_X}\left( E\times Y \right)$. Define
  \begin{align*}
    \mu\left( A,y \right) &= \mu_y\left( A_y \right)
  \end{align*}
  for each $A\in \mathcal{A}$. Here, we are using the fact that the projection onto $X$ produces the $\sigma$-algebra $\mathcal{B}_X$, and take the regular conditional measures with respect to the projection $\pi_Y\colon X\times Y\to Y$. Since $A_y\in \mathcal{A}_X$, it follows that $\mu\left( \cdot,y \right)$ is a probability measure on $\mathcal{A}$ for each $y\in Y$. If we have $A = \left( E\times Y \right)\cap \left( X\times B \right)$, and use the fact that $\pi_X\left( \left( E\times Y \right)^{y} \right)$, then we may write
  \begin{align*}
    \mu\left( A \right) &= \int_{Y}^{} \mu\left( A^y,y \right)\:d\mu_Y(y).
  \end{align*}
\end{proof}
\section{The Disintegration Theorem}%

\section{The Ergodic Decomposition}%

\nocite{measure_theory_bogachev,einsiedler_and_ward}
\printbibliography
\end{document}
